\documentclass[12pt]{article}
\usepackage[utf8]{inputenc}
\usepackage[T1]{fontenc}
\usepackage{hyperref}
\title{West Bengal Sign Language Language Resources for K-12 Curriculum Development}
\author{Compiled by Claude Sonnet 5 (claude-sonnet-5) for Andr\'as Kornai}
\date{}
\begin{document}
\maketitle

\section{Introduction}
West Bengal Sign Language (WBSL; ISO 639-3 wbs) is the sign language of the Kolkata Deaf community. Search results describe it as distinct from the Delhi variety of Indian Sign Language and related to Bangla Sign Language. The Calcutta Deaf and Dumb School, founded in 1883, is described as the first Deaf school in Bengal. No official WBSL dictionary or curriculum was found.

\section{Primary Educational Resources}
\begin{itemize}
\item Deaf Child Worldwide, ``Transforming Deaf Education: Supporting the futures of deaf children in West Bengal''. It describes work with CINI in Kolkata and notes that teachers often lack formal training in deafness: \url{https://www.ndcs.org.uk/deaf-child-worldwide/our-blog/transforming-deaf-education-supporting-the-futures-of-deaf-children-in-west-bengal}
\item NIOS Indian Sign Language course material, chapter 2 (ISL in general, not WBSL specifically): \url{https://www.nios.ac.in/media/documents/230-ISL/Ch-2.pdf}
\end{itemize}

\section{Linguistic and Reference Materials}
\begin{itemize}
\item Wikisigns West Bengal Sign Language, a user-created video dictionary of about 170 signs: \url{https://www.sign-lang.uni-hamburg.de/lr/compendium/lex/wikisignswestbengalsl.html}
\item Sign Language Studies 16(4), Gallaudet University Press: \url{https://ssl.gallaudet.edu/gupress/SLS/SLS16-4.html}
\item Project MUSE article: \url{https://muse.jhu.edu/article/627362}
\item FASAL proceedings 2026 (Reed College): \url{https://www.reed.edu/linguistics/khan/assets/FASALproc2026.pdf}
\item Deafness in Bangladesh, Wikipedia (related Bangla Sign Language context): \url{https://en.wikipedia.org/wiki/Deafness_in_Bangladesh}
\end{itemize}
I did not check the content of the Gallaudet, Project MUSE or FASAL items beyond their titles, so their relevance to WBSL is unconfirmed.

\section{Religious and Community Resources}
Scripture Earth lists a page for wbs: \url{https://scriptureearth.org/00eng.php?iso=wbs}. I did not verify what it contains.

\section{Digital and Audio Resources}
The Wikisigns dictionary above is the only video resource found.

\section{Curriculum Development Context}
West Bengal's Department of Mass Education and Extension and Netaji Subhas Open University ran an ISL workshop in 2025. The report is at \url{https://soe.wbnsou.ac.in/seminar/reports/20250922_Report_Workshop_on_ISL.pdf}. State training appears to focus on ISL rather than WBSL. A WBSL-medium curriculum would need documentation with the Deaf community and trained Deaf teachers.
\end{document}